\title{FlashAttention-3: Fast and Accurate Attention with Asynchrony and Low-precision}

\begin{abstract}
Serving as the foundational mechanism within the widely adopted Transformer architecture, attention constitutes a primary computational bottleneck in scaling large language models and managing long-context tasks. \textsc{FlashAttention} previously proposed an optimized attention implementation on GPUs by minimizing unnecessary memory traffic. Despite these improvements, recent hardware advancements have not been fully utilized; notably, \textsc{FlashAttention-2} leverages only 35\% of the H100 GPU’s capability. In this work, we present three key optimizations targeting Hopper GPUs: harnessing the asynchrony of Tensor Cores with TMA by (1) employing warp-specialization to maximize concurrency between computation and data transfer, (2) integrating block-wise matrix multiplication with softmax calculations, and (3) applying block quantization together with incoherent execution patterns, all supported by hardware-accelerated FP8 formats. Our solution, termed \textsc{FlashAttention-3}, achieves performance gains on H100 GPUs ranging from 1.5$\times$ to 2.0$\times$ over prior approaches, realizing up to 740 TFLOPs/s (a 75\% utilization) with FP16 precision and nearly 1.2 PFLOPs/s with FP8. Empirically, FP8 \textsc{FlashAttention-3} produces a 2.6$\times$ reduction in numerical error compared to a standard FP8-based attention mechanism.
\end{abstract}

\section{Introduction}
\label{sec:intro}

Within the Transformer architecture~\citep{vaswani2017attention}, the computational bottleneck predominantly arises from the attention mechanism, as calculating self-attention scores between queries and keys involves computation that grows quadratically with sequence length. Enhancing the scalability of attention mechanisms to accommodate longer contexts would enable new functionalities, such as modeling and reasoning across several extensive documents~\citep{guo2021longt5,shaham2022scrolls,peng2023yarn}, managing large codebases~\citep{roziere2023code, li2023starcoder}, as well as supporting diverse modalities including high-resolution visual content~\citep{chen2022scaling}, audio~\citep{gulati2020conformer}, and video~\citep{ho2022video}. Additionally, this expansion opens possibilities for new applications like maintaining extended histories in user interactions~\citep{sun2019bert4rec}, or coordinating agent operations over longer timeframes~\citep{yao2022react}. Accordingly, substantial attention has been devoted to accelerating attention mechanisms for long-context scenarios, either through approximate methods~\citep{katharopoulos2020transformers,choromanski2020rethinking,
tay2020efficient}, improvements in software efficiency~\citep{rabe2021self,dao2022flashattention,kwon2023efficient}, or by introducing alternative architectural approaches~\citep{peng2023rwkv,sun2023retentive,gu2023mamba}.

This paper extends the contributions of \citet{dao2022flashattention}, who introduced precise-attention methods informed by GPU architecture and execution patterns. In their work, \textsc{FlashAttention} was proposed—a new method based on tiling to facilitate efficient parallel processing of attention by merging all steps of attention into a single GPU kernel, thus avoiding slow global memory access for intermediate data. Subsequently, \citet{dao2023flashattention2} redesigned the algorithm as \textsc{FlashAttention-2}, enabling parallelism across the sequence length and optimizing the forward pass inner loop by segmenting it over blocks from the key and value matrices. These improvements enhanced workload balancing and resource utilization on GPUs. Nevertheless, our analysis indicates that \textsc{FlashAttention-2} underperforms in terms of resource usage on the latest GPUs when compared to highly optimized matrix-multiplication (GEMM) kernels, for example, achieving just 35\% utilization versus the 80-90\% seen with GEMM kernels on Hopper H100. This discrepancy is partly due to choices at the implementation level, such as employing Ampere instructions on Tensor Cores rather than taking advantage of Hopper-specific instructions. Prior research, including ThunkerKitten~\citep{spector2024thunder} and cuDNN 9~\citep{cudnn9}, demonstrates that leveraging Hopper-targeted instructions and a tiling-centric approach can accelerate attention mechanisms and streamline development.

At a foundational level, the algorithm implemented in \textsc{FlashAttention-2} operates within a basic synchronous paradigm and intentionally omits asynchronous mechanisms or low-precision considerations in its construction. Hardware-induced asynchrony emerges as a consequence of specialized modules aimed at expediting critical operations in machine learning pipelines, such as dedicated units for matrix multiplications (Tensor Cores) and distinct components for memory transfer (Tensor Memory Accelerator -- TMA), apart from the general-purpose CUDA cores managing logic, integer, and floating-point tasks. Lower numerical precisions—FP8 with Hopper and FP4 with Blackwell, continuing the progression from FP16 introduced in Pascal (2017) and BF16 in Ampere (2020)—consistently demonstrate increased computational throughput per watt and silicon real estate. The expanded functionalities available with Hopper are discussed in detail in \cref{subsec:hardware}. The central technical task thus involves refactoring \textsc{FlashAttention-2} to exploit these hardware advancements: asynchrony necessitates concurrent execution of matmul and softmax despite dependency constraints, while deploying low-precision arithmetic calls for strategic handling to mitigate quantization artifacts, particularly in the presence of rare or extreme values in LLMs~\citep{dettmers2208llm, sun2024massive}.

In pursuit of this goal, we introduce \textsc{FlashAttention-3}, presenting a synthesis of three innovative approaches designed to optimize performance on modern GPU platforms.\footnote{Although our experiments focus primarily on NVIDIA's Hopper architecture, the methodology applies broadly to any GPU architecture that supports strong asynchronous execution and low-precision computation.}

\begin{enumerate}

\item \textbf{Producer-Consumer asynchrony:} We introduce a warp-specialized software pipelining approach that leverages the asynchronous operation of data transfers and Tensor Core computations by assigning data producers and consumers to distinct warps. This separation increases the algorithm’s capacity to obscure both memory access and instruction scheduling delays.

\item \textbf{Hiding softmax under asynchronous block-wise GEMMs:} To mask the relatively low throughput of non-GEMM components within softmax—such as fused multiply-add and exponential calculations—we interleave their execution with asynchronous block-level GEMM instructions via WGMMA. In doing so, we restructure the \textsc{FlashAttention-2} algorithm to break certain inherent sequential constraints between softmax and GEMM operations. For instance, in the two-phase algorithm variant, while a block of the score matrix undergoes softmax processing, WGMMA is asynchronously launched to prepare the subsequent block.

\item \textbf{Hardware-accelerated low-precision GEMM:} We modify the forward computation to enable the use of FP8 Tensor Cores for GEMM operations, resulting in a near twofold improvement in observed TFLOPs/s. Achieving this entails reconciling the memory layout requirements between WGMMA—which expects blocks of FP32 accumulators and FP8 operands in specific arrangements. By employing block quantization together with incoherent block handling, we reduce the accuracy degradation typically caused by shifting to FP8 precision.
\end{enumerate}

For empirical evaluation, we assess the performance of \textsc{FlashAttention-3} on an H100 SXM5 GPU across a range of configurations. Our findings indicate that (1) in the forward pass, FP16 delivers 1.5-2.0$\times$ higher throughput compared to \textsc{FlashAttention-2}, achieving up to 740 TFLOPs/s, and realizes 1.5-1.75$\times$ faster computations during the backward pass; (2) FP8 attains nearly 1.2 PFLOPs/s; and (3) at larger sequence lengths, FP16 shows superior performance, with FP8 performing competitively\footnote{Specifically, FP8 outperforms for a head dimension of 64, while for head dimensions of 128 and 256, FP8 matches performance in scenarios without causal masking but falls behind in those with causal masking.} relative to NVIDIA’s cuDNN’s cutting-edge attention implementation. Additionally, we establish that FP16 \textsc{FlashAttention-3} retains the same numerical accuracy as \textsc{FlashAttention-2} and surpasses conventional attention implementations, since intermediate values (such as softmax scaling) are maintained in FP32. Furthermore, in contexts with outlier features, FP8 \textsc{FlashAttention-3} employing block quantization and incoherent processing exhibits 2.6$\times$ superior accuracy over standard attention approaches that use per-tensor quantization.

\textsc{FlashAttention-3} has been released as open-source under a permissive license\footnote{\textsc{FlashAttention-3}
  is available at \texttt{https://github.com/Dao-AILab/flash-attention}}. Furthermore, we intend to incorporate it into both the PyTorch and Hugging Face ecosystems in order to maximize its accessibility and impact for researchers and practitioners.

\section{Background: Multi-Head Attention and GPU Characteristics}
\label{sec:background}

\subsection{Multi-Head Attention}
\label{subsec:multi_head_attn}

Let $\mathbf{Q}, \mathbf{K}, \mathbf{V} \in \mathbb{R}^{N \times d}$ be the query, key and value input sequences associated to a single head, where $N$ is the sequence length and $d$ is the head dimension. Then the attention output $\mathbf{O}$ is computed as:

\begin{equation*}
  \mathbf{S} = \alpha \mathbf{Q} \mathbf{K}^\top \in \mathbb{R}^{N \times N}, \quad \mathbf{P} = \mathrm{softmax}(\mathbf{S}) \in \mathbb{R}^{N \times N}, \quad \mathbf{O} = \mathbf{P}\mathbf{V} \in \mathbb{R}^{N \times d},
\end{equation*}

In this context, the $\mathrm{softmax}$ function operates over each row, and it is common to choose $\alpha = 1/\sqrt{d}$ as the scale parameter. To mitigate numerical issues associated with the exponentiation step, $\mathrm{rowmax}(\mathbf{S})$ is subtracted from $\mathbf{S}$ as a stabilizing precaution. For multi-head attention (MHA), separate query, key, and value projection layers are defined for each attention head; the computation is executed in parallel for all heads and batches, resulting in the complete output tensor.

Now let $\phi$ be a scalar loss function and let $\mathbf{d}(-) = \partial \phi / \partial (-)$ be notation for the gradient. Given the output gradient $\mathbf{dO} \in \mathbb{R}^{N \times d}$, we compute $\mathbf{dQ}$, $\mathbf{dK}$, and $\mathbf{dV}$ according to the chain rule as follows:

\begin{align*}
  \mathbf{dV} &= \mathbf{P}^\top \mathbf{dO} \in \mathbb{R}^{N \times d} \\
  \mathbf{dP} &= \mathbf{dO} \mathbf{V}^\top \in \mathbb{R}^{N \times N} \\
  \mathbf{dS} &= \mathrm{dsoftmax} (\mathbf{dP}) \in \mathbb{R}^{N \times N} \\
  \mathbf{dQ} &= \alpha \mathbf{dS} \mathbf{K} \in \mathbb{R}^{N \times d} \\
  \mathbf{dK} &= \alpha \mathbf{dS}^\top \mathbf{Q} \in \mathbb{R}^{N \times d},
\end{align*}

In this context, the relationship $\mathbf{d}s = (\mathrm{diag}(p) - p p^\top)\mathbf{d}p$ holds when $p = \mathrm{softmax}(s)$, with $s$ representing a vector. The notation $\mathrm{dsoftmax}(\mathbf{dP})$ indicates the row-wise application of this expression. Moreover, this calculation is amenable to parallel execution over the heads and batches during the backward propagation stage in MHA.

\subsection{GPU hardware characteristics and execution model}
\label{subsec:hardware}

We outline the key features of the GPU execution model pertinent to \textsc{FlashAttention-3}, concentrating specifically on the NVIDIA Hopper architecture as a representative implementation of this model.

\paragraph{Memory hierarchy:} The GPU architecture consists of a multilayered arrangement of memory types, where greater bandwidth is associated with smaller storage capacities (\cref{tab:gpu-hierarchy})\footnote{\citet{luo2024benchmarking} provides a shared memory bandwidth value of 128 bytes per clock cycle per SM; this figure is scaled by the presence of 132 SMs and a boost frequency of 1830 MHz.}. At the top level, global memory (GMEM), or HBM, is located off-chip and is accessible by all streaming multiprocessors (SMs). Transfers from GMEM are automatically stored in an intermediate on-chip L2 cache. Each SM also possesses its own local, highly banked, programmer-managed shared memory (SMEM) of limited size. At the most localized level, SMs contain a register file for immediate data storage.

\paragraph{Thread hierarchy:} GPUs utilize a structured programming model that groups execution entities known as threads into a multi-level hierarchy. Starting from the most granular, the hierarchy includes individual threads, which are bundled into warps (each consisting of 32 threads), followed by warpgroups (groups of 4 consecutive warps). Above these are threadblocks—also referred to as cooperative thread arrays (CTAs)—then, in architectures like Hopper, threadblock clusters, and finally, grids at the highest organizational level.

The two hierarchies exhibit significant interdependence. Threads belonging to a single CTA are executed together on a shared SM, while CTAs grouped within a cluster are concurrently assigned to the same GPC. All threads inside a CTA have direct access to SMEM, but each individual thread possesses a maximum of 256 exclusive registers (RMEM).

\paragraph{Asynchrony and warp-specialization:}

GPUs function as throughput-oriented processors, leveraging high degrees of concurrency and asynchronous operations to mask the latencies associated with memory accesses and instruction execution. For performing asynchronous data transfers between GMEM and SMEM, Hopper introduces a specialized hardware component known as the Tensor Memory Accelerator (TMA) \cite[\S7.29]{cuda}. Additionally, diverging from older architectures like Ampere, Hopper’s Tensor Core—which can be accessed via the warpgroup-wide WGMMA instruction \cite[\S9.7.14]{ptx}—operates asynchronously and is capable of fetching its input data directly from shared memory.

Enabling hardware-level asynchrony permits the use of warp-specialized kernels, wherein the warps within a CTA are assigned exclusively as producers or consumers, performing only data movement or computation tasks, respectively. Such specialization generally enhances the compiler’s capacity to produce more efficient instruction schedules \citep{warp-specialization-2011}. Furthermore, Hopper introduces the ability to dynamically adjust register allocation among warpgroups using \verb|setmaxnreg| \cite[\S9.7.17.1]{ptx}, thus granting MMAs warps access to a larger portion of RMEM, while warps responsible solely for issuing TMA—requiring merely a single thread—are allocated fewer registers.

\paragraph{Low-precision number formats:}
\label{sec:low-precision-gpu}
Contemporary GPUs feature dedicated hardware designed to optimize calculations with reduced precision. As an illustration, the WGMMA instruction is capable of leveraging FP8 Tensor Cores within Hopper GPUs, which results in double the throughput per SM relative to FP16 or BF16 formats.

To effectively utilize FP8 WGMMA, one must pay careful attention to the operand layout requirements. Consider a GEMM operation computing $A \times B^{\top}$, where $A$ is an $M\times K$ matrix and $B$ is an $N\times K$ matrix. The operand is described as \emph{mn-major} if its memory layout is contiguous along the $M$ or $N$ dimensions, and as \emph{k-major} if contiguity is along the $K$ dimension instead. With FP16 WGMMA, the shared memory (SMEM) input operands can be supplied in either mn-major or k-major configuration. In contrast, FP8 WGMMA only accepts the k-major arrangement. Furthermore, in compound scenarios—such as attention mechanisms—where sequential GEMM operations are merged within one kernel, the mismatch between FP32 accumulator layouts and FP8 operand layouts presents a challenge for chaining FP8 WGMMA computations.

With regard to attention mechanisms, these constraints on layout necessitate specific adjustments to the structure of an FP8 algorithm, details of which are provided in \cref{sec:algofp8}.

\subsection{Standard Attention and Flash Attention} In accordance with \citet{dao2022flashattention}, we define \textbf{standard attention} as the GPU-based attention mechanism that stores the intermediate matrices $\mathbf{S}$ and $\mathbf{P}$ in HBM. \textsc{FlashAttention} was introduced to eliminate the costly memory operations associated with these intermediates by utilizing a localized softmax reduction and combining the attention computation into a unified kernel. This localized softmax process encompasses lines \ref{code-ws:softmax_start}-\ref{code-ws:softmax_end} in the consumer mainloop described in \cref{alg:flash3_wgmma_ws_only}, as well as the necessary rescaling of blocks in $\mathbf{O}$. A straightforward derivation demonstrating that this method accurately produces $\mathbf{O}$ is presented in \cite[\S 2.3.1]{dao2023flashattention2}.

\section{FlashAttention-3: Algorithm}
\label{sec:algo}

This section presents the \textsc{FlashAttention-3} algorithm, concentrating on the forward pass for clarity; details of the backward pass are provided in~\cref{sec:algo_ws_bwd}. We begin by demonstrating how warp-specialization can be combined with a circular SMEM buffer within the foundational structure of \textsc{FlashAttention-2}. Subsequently, we illustrate the use of WGMMA's asynchrony to implement a two-stage GEMM-softmax pipeline with overlapping execution. Lastly, we outline the adjustments necessary for FP8, covering aspects such as layout adherence and accuracy improvements through block quantization and non-coherent computation.

\subsection{Producer-Consumer asynchrony through warp-specialization and pingpong scheduling}
\label{sec:algo_ws}

\paragraph{Warp-specialization}
Similar to the approach taken in \textsc{FlashAttention-2}, the forward pass of \textsc{FlashAttention-3} exhibits high degrees of parallelism with respect to batch size, number of heads, and length of the query sequence. Therefore, an analysis at the CTA level suffices, wherein the algorithm processes a submatrix $\mathbf{Q}_i$ from the query to generate the associated output tile $\mathbf{O}_i$. For clarity, this discussion focuses initially on the warp-specialization design using a circular SMEM buffer, without considering the simultaneous GEMM-softmax operation. Let us denote $d$ as the dimensionality of each head, $N$ as the sequence length, and choose a query block size $B_r$ such that the query matrix $\mathbf{Q}$ is partitioned into $T_r = \lceil \frac{N}{B_r} \rceil$ segments $\mathbf{Q}_1, .., \mathbf{Q}_{T_r}$.

\begin{algorithm}[H]
    \caption{\small\label{alg:flash3_wgmma_ws_only}\textsc{FlashAttention-3} forward pass \textbf{excluding} intra-consumer overlap -- CTA perspective}
    \begin{algorithmic}[1]
\REQUIRE Input matrices $\mathbf{Q}_i \in \mathbb{R}^{B_r \times d}$ and $\mathbf{K}, \mathbf{V} \in \mathbb{R}^{N \times d}$ residing in HBM; select key block size $B_c$ and compute $T_c = \lceil \frac{N}{B_c} \rceil$.
\STATE Set up a pipeline manager to coordinate barrier sync, utilizing a circular SMEM buffer with $s$ stages.
\IF {warpgroup is producer}
\STATE Release a designated portion of registers.
\STATE Initiate transfer of $\mathbf{Q}_i$ from HBM to shared memory.
\STATE Once transfer completes, signal to consumers that $\mathbf{Q}_i$ is ready.
\FOR{$0 \le j < T_c$}
    \STATE Wait until consumers have finished using the $(j\,\%\,s)$ buffer stage.
    \STATE Trigger loads of $\mathbf{K}_j$ and $\mathbf{V}_j$ from HBM into shared memory at buffer position $(j\,\%\,s)$.
    \STATE After load, signal consumers of the ready state of $\mathbf{K}_j$ and $\mathbf{V}_j$.
\ENDFOR
\ELSE 
\STATE Allocate the predetermined register count based on consumer warp count.
\STATE On-chip, set $\mathbf{O}_i = (0) \in \mathbb{R}^{B_r \times d}$; initialize $\ell_i = 0$, $m_i = -\infty$ in $\mathbb{R}^{B_r}$.
\STATE Await the load completion of $\mathbf{Q}_i$ in shared memory.
\FOR{$0 \le j < T_c$}
\STATE Hold until $\mathbf{K}_j$ becomes available in shared memory.
\STATE Calculate $\mathbf{S}_i^{(j)} = \mathbf{Q}_i \mathbf{K}_j^T$ (using SS-GEMM). Perform commit and synchronize. \label{alg:ws_only_gemm1}
\STATE Save $m_i^{\mathrm{old}} = m_i$; update $m_i = \mathrm{max}(m_i^{\mathrm{old}}, \mathrm{rowmax}(\mathbf{S}_i^{(j)}))$. \label{code-ws:softmax_start}
\STATE Set $\widetilde{\mathbf{P}}_i^{(j)} = \mathrm{exp}(\mathbf{S}_i^{(j)} - m_i)$; update $\ell_i = \mathrm{exp}(m_i^{\mathrm{old}} - m_i) \ell_i + \mathrm{rowsum}(\widetilde{\mathbf{P}}_i^{(j)})$. \label{code-ws:softmax_end}
\STATE Wait for availability of $\mathbf{V}_j$ in shared memory.
\STATE Update $\mathbf{O}_i = \mathrm{diag}(\mathrm{exp}(m_i^{\mathrm{old}} - m_i))^{-1} \mathbf{O}_i + \widetilde{\mathbf{P}}_i^{(j)} \mathbf{V}_j$ via RS-GEMM. Commit and synchronize. \label{alg:ws_only_gemm2}
\STATE Release buffer stage $(j\,\%\,s)$ back to the producer.
\ENDFOR
\STATE Finalize $\mathbf{O}_i = \mathrm{diag}(\ell_i)^{-1} \mathbf{O}_i$ and $L_i = m_i + \log(\ell_i)$.
\STATE Store $\mathbf{O}_i$ and $L_i$ in HBM, corresponding to the $i$th output block of $\mathbf{O}$ and $L$.
\ENDIF
\end{algorithmic}
\end{algorithm}

In our adaptation of \cref{alg:flash3_wgmma_ws_only} for Hopper, we employ \verb|setmaxnreg| to manage (de)allocations, utilize TMA to load $\mathbf{Q}_i$ alongside $\{ \mathbf{K}_j, \mathbf{V}_j \}_{0 \leq j < T_c}$, and rely on WGMMA for executing GEMMs within the consumer mainloop. Here, the SS or RS designation specifies whether the first GEMM operand originates from shared memory or the register file, respectively. When analyzing the runtime behavior of \cref{alg:flash3_wgmma_ws_only}, it is important to recognize that TMA load commands are issued asynchronously and do not block waiting for other loads to finish. Additionally, in the producer mainloop, during the initial $s$ passes as the buffer is being populated, waits will not be triggered.

\paragraph{Pingpong scheduling} The asynchronous execution capabilities of WGMMA and TMA, combined with warp specialization, create possibilities for interleaving softmax computation from one warpgroup with GEMM tasks from another warpgroup. This approach is particularly compelling given that, on contemporary hardware accelerators, throughput for non-matmul operations lags significantly behind matmul throughput. For instance, the H100 SXM5 GPU achieves 989 TFLOPS for FP16 matmul operations but manages only 3.9 TFLOPS for special functions such as exponential\footnote{According to the CUDA programming guide, a single streaming multiprocessor (SM) can execute 16 special function operations per clock cycle. Therefore, multiplying 16 by 132 SMs and the 1830 MHz clock speed yields a peak of 3.9 TFLOPS for special functions.} (essential for softmax). When running the attention forward pass in FP16 with a head dimension of 128, matmul operations account for 512 times more FLOPS than exponential computations, whereas the exponential’s throughput is 256 times lower, leading to exponential calculations consuming about 50\% as many cycles as matmul. This discrepancy becomes even starker with FP8, since matmul throughput doubles without an increase in exponential throughput.

Because the computation of the exponential function is handled by a dedicated hardware component (the multi-function unit), it is optimal to overlap this operation with the execution of matrix multiplications on the Tensor Cores. To achieve this concurrency, we introduce synchronization barriers (\texttt{bar.sync} instructions), which enforce that the GEMMs (specifically, GEMM1 -- $\mathbf{P} \mathbf{V}$ from the current iteration, and GEMM0 -- $\mathbf{Q} \mathbf{K}^\top$ from the subsequent iteration) assigned to warpgroup 1 are executed prior to those assigned to warpgroup 2. Consequently, warpgroup 1 can carry out the softmax computation concurrently while warpgroup 2 is busy with its GEMMs. The warpgroup responsibilities then alternate, enabling warpgroup 2 to perform softmax while warpgroup 1 returns to GEMMs—an approach referred to as "pingpong" scheduling, visualized in~\cref{fig:pingpong_scheduling}. Although this scheduling pattern is not always as regular in practice as suggested by the diagram, empirical results generally show improved throughput (for example, increasing from 570 TFLOPS to between 620 and 640 TFLOPS in FP16 forward passes when using a head dimension of 128 and a sequence length of 8192).

\paragraph{Attention variants} In the case of multi-query attention~\citep{shazeer2019fast} and grouped query attention~\citep{ainslie2023gqa}, our implementation mirrors that of \textsc{FlashAttention-2}, modifying tensor indexing strategies to eliminate unnecessary duplication of $\mathbf{K}$ and $\mathbf{V}$ within HBM.

\subsection{Intra-warpgroup overlapping GEMMs and softmax}

Within a single warpgroup, it is possible to interleave certain instructions from the softmax operation with others from the GEMMs. Here, we outline a method to achieve this overlap.

In the attention algorithm, operations within the inner loop (main loop) have sequential dependencies that impede parallelization within a single iteration. For example, (local) softmax (lines \ref{code-ws:softmax_start} to \ref{code-ws:softmax_end}) relies on the output $\mathbf{S}_i^{(j)}$ of the first GEMM, while the second GEMM takes its result $\widetilde{\mathbf{P}}_i^{(j)}$ as an operand. Indeed, the wait statements in lines \ref{alg:ws_only_gemm1} and \ref{alg:ws_only_gemm2} of \cref{alg:flash3_wgmma_ws_only} serialize the execution of softmax and GEMMs. However, we can break these dependencies by pipelining across iterations through additional buffers in registers. Pursuing this idea, we propose the following two-stage\footnote{Note that the number of stages of the overlapping scheme is bounded by, but need not equal, the number $s$ of stages in the circular SMEM buffer.} GEMM-softmax pipelining algorithm:

\begin{algorithm}[H]
    \caption{\small\label{alg:flash3_wgmma}\textsc{FlashAttention-3} consumer warpgroup forward pass}
    \begin{algorithmic}[1]
      \REQUIRE  Matrices $\mathbf{Q}_i \in \mathbb{R}^{B_r \times d}$ and $\mathbf{K}, \mathbf{V} \in \mathbb{R}^{N \times d}$ located in HBM, key block size $B_c$, number of blocks $T_c = \lceil \frac{N}{B_c} \rceil$.
      \STATE Adjust the allocation of registers dynamically based on the total number of consumer warps.
      \STATE Set up on-chip variables: initialize $\mathbf{O}_i = (0) \in \mathbb{R}^{B_r \times d}$, $\ell_i = (0) \in \mathbb{R}^{B_r}$, and $m_i = (-\infty) \in \mathbb{R}^{B_r}$.
      \STATE Await availability of $\mathbf{Q}_i$ and $\mathbf{K}_0$ in shared memory.
      \STATE Use WGMMA to calculate $\mathbf{S}_{\mathrm{cur}} = \mathbf{Q}_i \mathbf{K}_0^T$. Commit and block until completion.
      \STATE Free buffer segment corresponding to the initial stage ($0$) for $\mathbf{K}$.
      \STATE Determine $m_i$, $\tilde{\mathbf{P}}_{\mathrm{cur}}$, and $\ell_i$ using $\mathbf{S}_{\mathrm{cur}}$, then update the scale of $\mathbf{O}_i$.
      \FOR{$1 \leq j < T_c - 1$}
        \STATE Wait until $\mathbf{K}_j$ is available in shared memory.
        \label{code:mainloop_start}
        \STATE Perform WGMMA to compute $\mathbf{S}_{\mathrm{next}} = \mathbf{Q}_i \mathbf{K}_j^T$; commit without awaiting completion. \label{code:first_wgmma}
        \STATE Wait for $\mathbf{V}_{j-1}$ to become accessible in shared memory.
        \STATE Execute WGMMA for $\mathbf{O}_i = \mathbf{O}_i + \tilde{\mathbf{P}}_{\mathrm{cur}} \mathbf{V}_{j-1}$; commit without blocking. \label{code:second_wgmma}
        \STATE Wait for the completion of the WGMMA for $\mathbf{Q}_i \mathbf{K}_j^T$.
        \STATE Update $m_i$, $\tilde{\mathbf{P}}_{\mathrm{next}}$, and $\ell_i$ based on $\mathbf{S}_{\mathrm{next}}$. \label{code:softmax}
        \STATE Await the finish of WGMMA for $\tilde{\mathbf{P}}_{\mathrm{cur}} \mathbf{V}_{j-1}$, and then rescale $\mathbf{O}_i$.
        \STATE Release the buffer stage for $\mathbf{K}$ at $(j\,\%\,s)$, and for $\mathbf{V}$ at $(j-1\,\%\,s)$.
        \STATE Transfer $\mathbf{S}_{\mathrm{next}}$ contents into $\mathbf{S}_{\mathrm{cur}}$.
        \label{code:mainloop_end}
      \ENDFOR \STATE Wait until $\mathbf{V}_{T_c - 1}$ is resident in shared memory.
      \STATE Use WGMMA to evaluate $\mathbf{O}_i = \mathbf{O}_i + \tilde{\mathbf{P}}_{\mathrm{last}} \mathbf{V}_{T_c - 1}$, committing and waiting for completion.

\STATE Epilogue: Adjust $\mathbf{O}_{i}$ according to $m_{i}$. Determine $L_{i}$ utilizing $m_{i}$ and $\ell_{i}$. 
Store $\mathbf{O}_{i}$ and $L_{i}$ in HBM, assigning them as the $i$-th segments of $\mathbf{O}$ and $L$. 
\end{algorithmic}
\end{algorithm}

\cref{alg:flash3_wgmma} substitutes for the consumer pathway described in \cref{alg:flash3_wgmma_ws_only}, thereby forming the full \textsc{FlashAttention-3} procedure for FP16 computations. Conceptually, WGMMA is utilized as shorthand for asynchronous GEMM. Throughout the mainloop (spanning lines \ref{code:mainloop_start} to \ref{code:mainloop_end}), the subsequent WGMMA call in iteration $j$ (line \ref{code:second_wgmma}) proceeds concurrently with the execution of softmax calculations belonging to iteration $j+1$ (line \ref{code:softmax}).

Although the pipelined architecture depicted earlier is theoretically advantageous for performance, several practical concerns arise:
\paragraph{Compiler reordering} The ideal sequence shown in the pseudocode may not be maintained during compilation, as NVCC frequently rearranges instructions to optimize execution. This process can disrupt the intended interleaving of WGMMA and non-WGMMA operations in the pipeline, resulting in less predictable execution patterns or reduced performance improvements. Examination of SASS reveals that the compiler does produce overlapped instruction streams as intended (see Section~\ref{sec:2-stage-sass}).
\paragraph{Register pressure} Achieving efficient pipelining relies on minimizing register spilling. The 2-stage pipeline, however, imposes extra register requirements because it must hold intermediate values and context across pipeline stages. In particular, each threadblock is required to store an additional $\mathbf{S}_{\mathrm{next}}$, which incurs a supplementary register allocation of $B_r \times B_c \times \text{sizeof}(\text{float})$. This elevated register usage could be at odds with choosing larger block sizes, as both strategies compete for register resources. It is advisable to make configuration choices based on empirical profiling.
\paragraph{3-stage pipelining} Building upon the prior 2-stage design, a 3-stage strategy is proposed that enables the second WGMMA operation to be further overlapped with the softmax computation. This refinement has the potential to exploit Tensor Core resources more fully, but necessitates even greater register usage since it adds another phase to the pipeline. Consequently, balancing pipeline depth against tile size becomes more challenging. Full details and empirical assessments of the 3-stage algorithm are provided in~\cref{sec:3-stage}.

\subsection{Low-precision with FP8}
\label{sec:algofp8}

\textbf{Efficiency: layout transformations.} Executing the forward pass of \textsc{FlashAttention-3} using FP8 precision introduces unique difficulties related to layout compatibility, which do not arise when working with FP16.

To begin, we observe that the input tensors $\mathbf{Q}$, $\mathbf{K}$, and $\mathbf{V}$ are usually arranged so that their head dimension is contiguous. However, for the second GEMM operation utilizing FP8 WGMMA, the k-major requirement stipulates that $\mathbf{V}$—specifically, the portions of $\mathbf{V}$ transferred into SMEM—must instead be contiguous along the sequence length. Given that TMA loading cannot alter which dimension is contiguous, there are two available approaches: (1) pre-transpose $\mathbf{V}$ in GMEM before use, or (2) perform an in-kernel transpose on chunks of $\mathbf{V}$ after they have been placed in SMEM. To realize the first approach, one could (1a) integrate the transpose with the output stage of an earlier process, such as the rotary embedding, or (1b) employ an independent transpose kernel as a preparatory step\footnote{A well-tuned transpose kernel is capable of operating close to the hardware’s peak memory bandwidth~\citep{colfax_cutlass_transpose_2024}.} to interchange the order of the sequence length and head dimensions. Nonetheless, (1a) is challenging to incorporate in a standard library context, and (1b) imposes excessive overhead in scenarios dominated by memory limitations, such as inference workloads.

For FP8 \textsc{FlashAttention-3}, we select strategy (2). To implement the in-kernel transpose, we utilize LDSM (\verb|ldmatrix|) and STSM (\verb|stmatrix|) instructions, which coordinate a warp of threads to load data from shared memory (SMEM) to register memory (RMEM) and store from RMEM to SMEM in blocks of 128 bytes.\footnote{According to the PTX reference, LDSM/STSM are intended for moving $8 \times 8$ matrices with 16-bit elements \cite[\S 9.7.13.4.15-16]{ptx}. We overcome this limitation by packing pairs of 8-bit values, thereby adapting LDSM/STSM for FP8 formats. Nevertheless, transpose operations with LDSM/STSM cannot separate packed 8-bit elements, which compels us to perform specific register shuffles between the LDSM and STSM calls in order to realize a per-tile transpose; those technical specifics are not elaborated here.} These instructions make highly efficient use of registers, allowing execution within the producer warpgroup, while also supporting layout transposes concurrent with memory movement. Additionally, beyond the initial pass, we arrange the scheduling so that transposing the next $\mathbf{V}$ tile is overlapped with the shadow execution of the two WGMMAs corresponding to the previous $\mathbf{V}$ and the present $\mathbf{K}$ tile.

Second, it becomes evident that, unlike FP16, the FP32 accumulator layout utilized in FP8 WGMMA deviates from the arrangement expected for operand A when stored in registers. Illustrative fragments of these layouts are presented in \cref{fig:rmem_accum} and \cref{fig:rmem_operand}, highlighting the sequence in which each thread maintains register entries. Employing byte permute instructions enables conversion of the first WGMMA's accumulator into a configuration compatible with the subsequent WGMMA and aligns with the $\mathbf{V}$ tile layout produced by the in-kernel transpose. More precisely, as shown in \cref{fig:rmem_accum}, the register order is adjusted to
$$\{ \verb|d0 d1 d4 d5 d2 d3 d6 d7| \},$$
with this register rearrangement repeated for every 8-byte block. This operation essentially permutes the columns of the $\mathbf{P}$ tile—for instance, columns $0189$ now precede the others. To guarantee WGMMA generates the appropriate output tile, the in-kernel transpose can be designed to perform an analogous row permutation for the $\mathbf{V}$ tile.\footnote{By carrying out the transpose within the kernel, we gain extra flexibility, thereby removing the need for shuffle instructions to shift register ownership between threads, as previously discussed in~\citep{colfax_fp8_flashattention_2024}.}

\textbf{Accuracy: block quantization and incoherent processing.} The FP8 (e4m3) representation allocates just 3 bits to the mantissa and 4 bits to the exponent, which inherently introduces greater numerical inaccuracies compared to FP16/BF16. Additionally, in large-scale models, the presence of outlier values with substantially larger magnitudes than most elements~\citep{dettmers2208llm, sun2024massive} exacerbates the challenges of quantization. Standard practice involves per-tensor scaling~\citep{micikevicius2022fp8}, where a single scaling factor is maintained for each tensor (for example, one each for $\mathbf{Q}$, $\mathbf{K}$, and $\mathbf{V}$). To counteract the increased numerical errors associated with FP8 attention, we utilize two primary approaches:

\begin{enumerate}

\item \textbf{Block quantization}: a single scalar is retained for each block, so for the tensors $\mathbf{Q}$, $\mathbf{K}$, and $\mathbf{V}$, we partition them into blocks of dimensions $B_r \times d$ or $B_c \times d$ and quantize each block individually. This approach allows the quantization step to be integrated seamlessly with pre-attention operations such as rotary embedding, without incurring additional overhead, as these operations are constrained by memory bandwidth. Given that the \textsc{FlashAttention-3} algorithm processes data in block format, we can efficiently rescale each block of $\mathbf{S}$ to reflect block-level quantization without extra computational demand.

\item \textbf{Incoherent processing}: to mitigate the impact of outlier values, we pre-multiply $\mathbf{Q}$ and $\mathbf{K}$ by a random orthogonal matrix $\mathbf{M}$ prior to FP8 quantization. As $\mathbf{M}$ satisfies $\mathbf{M} \mathbf{M}^\top = I$, it follows that $(\mathbf{Q} \mathbf{M}) (\mathbf{K} \mathbf{M})^\top = \mathbf{Q} \mathbf{K}^\top$, indicating the attention results remain unchanged after the transformation. This process disperses outliers, as each component of $\mathbf{Q} \mathbf{M}$ or $\mathbf{K} \mathbf{M}$ consists of a randomized linear combination of the original entries, thereby minimizing quantization artifacts. Consistent with the methodologies of \citet{chee2024quip} and~\citet{tseng2024quip}, we select $\mathbf{M}$ as a product of random diagonal matrices with entries $\pm 1$ and a Hadamard matrix, enabling multiplication in $O(d \log d)$ time as opposed to $O(d^2)$. Moreover, this can be merged with the rotary embedding operation without introducing extra computation.

We demonstrate in \cref{sec:numerical_error} that these methods can decrease numerical error by as much as 2.6$\times$.

\section{Empirical Validation}
\label{sec:exp}

Our implementation of \textsc{FlashAttention-3} utilizes foundational components from CUTLASS~\citep{Thakkar_CUTLASS_2023}, specifically the WGMMA and TMA abstractions, to assess both its performance and precision.

\begin{itemize}

\item \textbf{Benchmarking attention.} We evaluate the execution time of \textsc{FlashAttention-3} at varying sequence lengths and benchmark its performance against the standard PyTorch implementation, \textsc{FlashAttention-2}, a Triton-based version of \textsc{FlashAttention-2} (leveraging H100-optimized instructions), as well as the cuDNN vendor-supplied \textsc{FlashAttention-2} tuned for H100 GPUs. Our experiments demonstrate that \textsc{FlashAttention-3} achieves speeds up to 2.0$\times$ that of \textsc{FlashAttention-2} and is 1.5$\times$ faster than the Triton variant. Notably, \textsc{FlashAttention-3} attains a throughput of up to 740 TFLOPs/s, representing 75\% of H100 GPUs' peak theoretical TFLOPs/s.
\item \textbf{Ablation study.} Our results establish that the algorithmic enhancements introduced—namely, warp specialization and the integration of GEMM-softmax pipelining—are key factors underpinning the speed improvements seen with \textsc{FlashAttention-3}.
\item \textbf{Accuracy of FP8 attention.} We verify that by using block quantization and incoherent processing, the numerical errors encountered with FP8 \textsc{FlashAttention-3} are decreased by 2.6$\times$.
\end{itemize}

\subsection{Benchmarking Attention}
\label{subsec:benchmark_attn}

We evaluate the runtime performance of various attention mechanisms on an H100 80GB SXM5 GPU across multiple configurations (with and without causal masking, head dimensions set to either 64 or 128), utilizing FP16 data type for the input. The corresponding measurements are presented in~\cref{fig:benchmark_attn_fwd} and~\cref{fig:benchmark_attn_bwd}. The results indicate that \textsc{FlashAttention-3} delivers approximately 1.5-2.0$\times$ acceleration relative to \textsc{FlashAttention-2} during the forward computation, and achieves a 1.5-1.75$\times$ speedup during the backward computation. In comparison to a conventional attention implementation, \textsc{FlashAttention-3} provides a substantial speed advantage, ranging from 3$\times$ to as much as 16$\times$. Notably, for sequence lengths of 1k and above, corresponding to medium and large sequences, \textsc{FlashAttention-3} is able to outperform even a proprietary, highly optimized vendor library (cuDNN) designed specifically for H100 GPUs.

\paragraph{Benchmark settings:} We vary the sequence length as 512, 1k, ..., 16k, and set batch size so that the total number of tokens is 16k. We set the hidden dimension to 2048, and head dimension to be either 64, 128, or 256 (i.e., 32 heads, 16 heads, or 8 heads). To calculate the FLOPs of the forward pass, we use:

\begin{equation*}
  4 \cdot \text{seqlen}^2 \cdot \text{head dimension} \cdot \text{number of heads}.
\end{equation*}

Using causal masking, we reduce this figure by half, as only about 50% of the elements are computed. To determine the FLOPs required for the backward pass, we scale the forward pass FLOPs by a factor of 2.5. This adjustment reflects that the backward pass consists of five matrix multiplications, compared to two in the forward pass, mainly as a result of recomputation.

Additionally, we evaluate the runtime performance of FP8 during the forward pass with comparable configurations. The outcomes for headdim 256 are presented in ~\cref{fig:benchmark_attn_fp8}, while comprehensive data can be found in \cref{sec:benchmark_attn_fp8_full}.

\subsection{Ablation Study: 2-Stage Pipelining Experiments}
\label{sec:2-stage-pipelining-experiments}

We conduct an ablation study on both the 2-stage WGMMA-softmax pipelining approach and the use of warp-specialization within the non-causal FP16 \textsc{FlashAttention-3} framework, holding the parameters fixed at $\{\text{batch}, \text{seqlen}, \text{nheads}, \text{hdim}\} = \{ 4, 8448, 16, 128\}$. As shown in~\cref{table:ablation_pipelining}, the results validate that introducing asynchrony through warp-specialization, as well as enabling overlap between GEMM and softmax, provides considerable performance gains—increasing throughput from 570 TFLOPs to 661 TFLOPs.

\subsection{Numerical Error Validation}
\label{sec:numerical_error}

Given increased scrutiny regarding the numerical error~\citep{golden2024flash} associated with
\textsc{FlashAttention}, we evaluate and contrast \textsc{FlashAttention-2}, \textsc{FlashAttention-3}, and a conventional attention implementation relative to a FP64 reference version. In order to mimic the occurrence of outlier features and activations commonly seen in LLMs~\citep{dettmers2208llm,
  sun2024massive}, we sample the elements of $\mathbf{Q}, \mathbf{K}, \mathbf{V}$ according to the following distribution:

\begin{equation*}
  \mathcal{N}(0, 1) + \mathcal{N}(0, 100) \cdot \mathrm{Bernoulli}(0.001).
\end{equation*}

Specifically, all entries follow a normal distribution with mean zero and a standard deviation of 1, except for 0.1\% of the entries where an additional independent normal component with standard deviation 10 is introduced. Subsequently, we compute the root mean squared error (RMSE) as displayed in~\cref{table:numerical_error}. In the FP16 setting, both \textsc{FlashAttention-2} and \textsc{FlashAttention-3} achieve an RMSE that is 1.7$\times$ lower than the conventional approach, since the intermediate (softmax) values are maintained in FP32 precision. For the FP8 baseline attention, per-tensor scaling is used, the matrix multiplication accumulates results in FP32, and the intermediate softmax values are preserved in FP16. Owing to the employment of block quantization and incoherent processing, \textsc{FlashAttention-3} operating in FP8 offers 2.6$\times$ better accuracy compared to this baseline.

\section{Dicussion, Limitations, Conclusion}
\label{sec:discussion}

By leveraging \textsc{FlashAttention-3}, we illustrate that adopting novel programming paradigms alongside advancements in hardware—specifically asynchrony and reduced-precision computations—significantly boosts both the speed and precision of attention mechanisms. This approach achieves a 1.5-2.0$\times$ improvement in throughput relative to \textsc{FlashAttention-2}, and mitigates FP8 quantization errors by a factor of 2.6$\times$ over typical per-tensor quantization schemes. Some areas requiring further investigation include refining our method for LLM inference, implementing a persistent kernel architecture within the FP8 kernel,\footnote{In our experiments, FP16 \textsc{FlashAttention-3} is equipped with a persistent kernel and a load-balancing mechanism, while FP8 \textsc{FlashAttention-3} currently lacks these features. This distinction partially accounts for the inferior performance of FP8 \textsc{FlashAttention-3} in scenarios involving short sequences and causal masking compared to FP8 cuDNN kernels.} and probing the implications of low-precision attention in the context of large-scale model training. While this study centers on Hopper GPUs, we anticipate that the methods presented will generalize to other accelerator platforms. The availability of a more efficient and accurate attention primitive stands to facilitate innovation in tasks requiring extended context.

\end{document}